\documentclass[11pt, a4paper]{article}

% ── Packages ──────────────────────────────────────────────────────────────────
\usepackage[a4paper, margin=2.5cm]{geometry}
\usepackage{booktabs}
\usepackage{array}
\usepackage{hyperref}
\usepackage{parskip}
\usepackage{longtable}
\usepackage{enumitem}
\usepackage{microtype}

\hypersetup{colorlinks=true, linkcolor=blue, urlcolor=blue}

% ══════════════════════════════════════════════════════════════════════════════
\begin{document}

\pagenumbering{roman}

% ── Title block ───────────────────────────────────────────────────────────────
\begin{center}
  {\large Data Mining BSS 2026}\\[1.5cm]
  {\LARGE\bfseries Predicting Game Success on the Steam Platform}\\[1cm]
  Team 9\\
  Brewed Clusters\\[0.3cm]
  Satwik Kumar Lohani\\
  Arpitha Shivaprasad Kori\\
  Zubair Ashfaq\\
  Beyza Ünlü\\
  Muhammad Sameer Siddiqui\\
  Esha Raheel\\[0.5cm]
  April 2026\\[2cm]
  \small Submitted to\\
  Data and Web Science Group\\
  University of Mannheim
\end{center}

\newpage

% ══════════════════════════════════════════════════════════════════════════════
\section{Introduction to the task}

The Steam marketplace hosts over 50{,}000 active titles, making it one of the most
competitive digital distribution platforms in the world. Each year, thousands of new
games are released, yet the majority receive little to no player attention. Game
developers and publishers routinely face high-stakes decisions about genre, pricing,
and release timing with limited quantitative guidance.

This raises a non-trivial prediction challenge: a game's reception is not simply
determined by its genre or price alone. The interaction between content type, community
tags, platform availability, and pricing creates a complex, high-dimensional signal
that is difficult to reason about intuitively. The goal of this project is to
\textbf{predict whether a Steam game will receive a predominantly positive or negative
review outcome} based on its observable metadata attributes such as genre, price,
tags, and platform availability.

Beyond prediction, we aim to produce interpretable insight --- identifying which
attributes are most strongly associated with positive reception using SHAP values
and feature importances, providing actionable guidance for developers making design
and pricing decisions.

% ══════════════════════════════════════════════════════════════════════════════
\section{What data will we use?}

\begin{itemize}
  \item \textbf{Dataset:} Steam Games Dataset
  \item \textbf{Source:} Kaggle --- FronkonGames \footnote{\url{https://www.kaggle.com/fronkongames/steam-games-dataset}}, last updated 2024
  \item \textbf{Data Retrieval:}
  \begin{itemize}
    \item The dataset contains metadata for approximately 122{,}611 Steam game entries
          retrieved via the Steam API and SteamSpy.
    \item Games with zero user reviews cannot be assigned a reliable label and will
          be excluded. We use the remaining games with at least one review for
          classification.
    \item The binary target label is derived from the review ratio:
          \texttt{review\_ratio} $=$ Positive $\div$ (Positive + Negative).
          A game is labelled \textbf{Good} if \texttt{review\_ratio} $\geq$ 0.70,
          and \textbf{Bad} otherwise. Class imbalance will be handled via SMOTE
          oversampling and \texttt{class\_weight=`balanced'}.
  \end{itemize}
  \item \textbf{Feature Selection:} Table~\ref{tab:features} describes the features
        used in our classification task.
\end{itemize}

\begin{center}
\begin{longtable}{@{} p{4.2cm} p{2.5cm} p{7.8cm} @{}}
\caption{Description of features used in the classification task}
\label{tab:features}\\
\toprule
\textbf{Feature} & \textbf{Type} & \textbf{Description} \\
\midrule
\endfirsthead
\toprule
\textbf{Feature} & \textbf{Type} & \textbf{Description} \\
\midrule
\endhead
\midrule
\multicolumn{3}{r}{\small\itshape continued on next page}\\
\endfoot
\bottomrule
\endlastfoot
Price                          & Numeric     & Listed price in USD; 0 for free-to-play titles \\
Release era (derived)          & Categorical & Release year binned into era buckets (pre-2015, 2015--2019, 2020+) \\
Genres                         & Multi-value & Steam genre labels (e.g., Action, Indie, RPG); binarised via \texttt{MultiLabelBinarizer} \\
Tags (top-50)                  & Multi-value & User-defined gameplay tags; top 50 by frequency retained as binary features \\
Categories                     & Multi-value & Feature labels (e.g., Single-player, Co-op, VR); binarised \\
Achievements                   & Numeric     & Number of in-game achievements \\
Average playtime               & Numeric     & Average total playtime across all owners (minutes) \\
Recommendations                & Numeric     & Number of Steam recommendations \\
Windows / Mac / Linux          & Boolean     & Platform availability flags \\
Number of languages (derived)  & Numeric     & Count of supported languages \\
Required age                   & Numeric     & Minimum age rating \\
DLC count                      & Numeric     & Number of downloadable content packs \\
\texttt{review\_ratio} (derived) & Numeric   & Positive $\div$ (Positive + Negative); used to construct the target \\
\texttt{label} (target)        & Binary      & \textbf{Good} if \texttt{review\_ratio} $\geq$ 0.70, else \textbf{Bad} \\
\end{longtable}
\end{center}

\noindent\textit{Note: Raw text fields (game name, developer, publisher) are not used
directly as features. Free-text carries little transferable signal across games and
would require NLP treatment beyond the scope of this project.}

% ══════════════════════════════════════════════════════════════════════════════
\section{How will we solve the problem?}

\subsection{Related Work}

Several public analyses on the same Kaggle dataset have attempted to predict Steam
game reception using review-based labels. Prior approaches have largely relied on
simple feature sets (price, genre counts) with Logistic Regression or Random Forest,
reporting accuracy as the primary metric --- which is misleading under class imbalance.
Our work improves on this by: (1) applying a corrected CSV loading pipeline that
fixes a known column-misalignment bug in the raw data, (2) engineering richer
features from multi-value tag and genre fields, and (3) evaluating under macro F1
and Precision-Recall AUC, which are more appropriate for imbalanced binary
classification.

\subsection{Data Preprocessing}

\begin{itemize}
  \item \textbf{Removing unlabelled games:} Games with zero total reviews are excluded,
        as no label can be derived for them.
  \item \textbf{Parsing multi-value fields:} Genre, tag, and category fields are
        stored as comma-separated strings and are exploded into binary indicator
        columns using \texttt{MultiLabelBinarizer}.
  \item \textbf{Handling missing values:} Median imputation is applied for missing
        numeric values. Columns with very high missing rates that carry no predictive
        signal are dropped entirely.
  \item \textbf{Normalising numerical data:} Price is log-transformed to reduce right
        skew caused by the prevalence of free-to-play titles. All continuous features
        are standardised using \texttt{StandardScaler}.
  \item \textbf{Handling class imbalance:} SMOTE oversampling and
        \texttt{class\_weight=`balanced'} are used to address the imbalance between
        well-received and poorly-received games.
  \item \textbf{Reproducibility:} The cleaned dataset is exported to Parquet format
        to ensure consistent input across all notebooks.
\end{itemize}

\subsection{Exploratory Data Analysis}

Before modelling, the following analyses are planned to understand the data and
guide preprocessing decisions:

\begin{itemize}
  \item Distribution of review ratios and visual assessment of class balance
  \item Price distribution stratified by free-to-play vs.\ paid games
  \item Frequency distribution of top genres and user-defined tags
  \item Correlation heatmap across all numeric features
  \item Outlier identification in review counts and playtime using the IQR method
\end{itemize}

\subsection{Algorithms}

\subsubsection{Classification Models (Predict Review Outcome)}

\begin{itemize}
  \item \textbf{Majority Class Classifier:} A na\"{i}ve baseline that always predicts
        the dominant class. Establishes a lower performance bound.
  \item \textbf{Logistic Regression (L2):} A linear and interpretable model whose
        coefficient signs reveal which features push a game toward a Good or Bad
        prediction.
  \item \textbf{Random Forest:} An ensemble model that captures non-linear feature
        interactions and provides built-in feature importances for interpretability.
  \item \textbf{XGBoost:} A gradient-boosted tree model that is robust to class
        imbalance and is expected to achieve the highest predictive performance.
        Hyperparameters are tuned via nested cross-validation.
\end{itemize}

\subsubsection{Feature Importance (Identify Key Success Factors)}

\begin{itemize}
  \item \textbf{Random Forest feature importances:} Global ranking of which features
        contribute most to classification decisions, based on node impurity reduction.
  \item \textbf{SHAP (SHapley Additive Explanations):} Applied to the XGBoost model
        to provide both global feature importance and local explanations showing why
        a specific game was classified as Good or Bad. This directly addresses the
        business question by identifying which attributes developers should focus on.
\end{itemize}

% ══════════════════════════════════════════════════════════════════════════════
\section{How will we measure the success?}

\begin{itemize}
  \item \textbf{Evaluation metrics and prioritisation}

  Primary Metric: Macro F1-Score
  \begin{itemize}
    \item Why Macro F1? It averages F1 equally across both classes, preventing the
          majority class from masking poor minority-class performance --- critical for
          imbalanced datasets.
    \item But: solely relying on F1 may overlook important precision-recall trade-offs.
  \end{itemize}

  \item \textbf{Additional Metrics}
  \begin{itemize}
    \item AUC-ROC: measures the model's ability to discriminate between classes
          across all decision thresholds.
    \item Precision-Recall AUC: particularly informative under class imbalance, as
          it focuses on performance on the positive (minority) class.
    \item Confusion matrices reported per model to inspect false positive and false
          negative trade-offs explicitly.
    \item All metrics reported as mean $\pm$ std across 10 stratified folds.
  \end{itemize}

  \item \textbf{Hyperparameter Optimisation}
  \begin{itemize}
    \item Nested cross-validation: inner 3-fold loop for hyperparameter search,
          outer 10-fold loop for unbiased final evaluation.
    \item Key parameters tuned: regularisation strength $C$ for Logistic Regression;
          number of estimators and maximum depth for Random Forest; learning rate,
          maximum depth, and subsampling ratio for XGBoost.
  \end{itemize}

  \item \textbf{State-of-the-Art Comparison}
  \begin{itemize}
    \item Results will be compared against 2--3 public Kaggle kernels on the same
          dataset using macro F1 and AUC-ROC as the basis for comparison.
  \end{itemize}
\end{itemize}

% ══════════════════════════════════════════════════════════════════════════════
\section{References}

\begin{enumerate}
  \item FronkonGames, \textit{Steam Games Dataset}, Kaggle, 2024.\\
        \url{https://www.kaggle.com/fronkongames/steam-games-dataset}
  \item Lundberg, S.\ M., \& Lee, S.\ I.\ (2017). A unified approach to interpreting
        model predictions. \textit{Advances in Neural Information Processing Systems}, 30.
  \item Chen, T., \& Guestrin, C.\ (2016). XGBoost: A scalable tree boosting system.
        \textit{Proceedings of KDD 2016}.
\end{enumerate}

% ══════════════════════════════════════════════════════════════════════════════
\end{document}
